\documentclass[12pt]{article}
\usepackage[utf8]{inputenc}
\usepackage[T1]{fontenc}
\usepackage{hyperref}
\title{Poyan\'awa Language Resources for K-12 Curriculum Development}
\author{Compiled by Claude Sonnet 5.5 (claude-sonnet-5-5) for Andr\'as Kornai}
\date{}
\begin{document}
\maketitle

\section{Introduction}
Poyan\'awa (Puyanawa; ISO 639-3 pyn; Glottocode poya1241) is a Panoan language of Acre, Brazil. Search results describe it as nearly extinct, with a single fluent speaker reported as of 2023 and limited results from revival efforts. The searches found no dictionary, grammar or school materials for it. Note that Nova Vida, which appeared in one search, is a Shanenawa community and not Poyan\'awa.

\section{Primary Educational Resources}
None found.

\section{Linguistic and Reference Materials}
\begin{itemize}
\item Wikipedia: Poyanawa language. \url{https://en.wikipedia.org/wiki/Poyanawa_language}
\item Article proposing an internal classification of the Pano languages (context only). \url{https://periodicos.ufpe.br/revistas/INV/article/view/1443}
\end{itemize}

\section{Religious and Community Resources}
\begin{itemize}
\item Povos Ind\'igenas no Brasil (Instituto Socioambiental), entry on the Puyanawa people. \url{https://pib.socioambiental.org/en/povo/Puyanawa}
\end{itemize}

\section{Digital and Audio Resources}
\begin{itemize}
\item Endangered Languages Project entry for Poyan\'awa. \url{https://endangeredlanguages.com/lang/2949}
\end{itemize}

\section{Curriculum Development Context}
With almost no documentation found and at most one speaker, any heritage curriculum would depend on community-led work and on whatever recordings and field notes exist, which this search did not locate. The PIB entry is the best starting point for community and contact context.

\end{document}